\chapter{Concluding remarks and further work}
\label{Ch:ConcludingRemarksFurtherWork}

The two papers in this thesis together give a calibration and validation
route for ductile-fracture prediction in aluminium sheet across the full
range of stress states seen in a crash event, rather than the
tension-dominated subset typically covered. The rate-dependent necking
criterion from Paper 1 and the triaxiality-dependent damage model from
Paper 2 share a common experimental basis, which is what allows the two to
be combined directly in a crash simulation without a separate
re-calibration step.

Further work should extend the compression-shear specimen and calibration
methodology to other automotive aluminium alloys, and evaluate whether the
rate-sensitivity parameter identified in Paper~1 needs to be re-calibrated
at the higher strain rates seen in a real crash pulse, which exceed the
$10^{2}$~s$^{-1}$ upper bound tested here.
